\begin{table}[H]
\centering\sffamily
\footnotesize
\setlength{\tabcolsep}{3pt}
\caption{Men without a recorded interval: bounds and inverse probability weighting}
\label{tab:s8}
\begin{tabular}{@{}lrrrrrr@{}}
\toprule
\rowcolor{tablehead}\makecell[l]{County rurality} & \makecell[r]{Treated\\men} & \makecell[r]{Recorded\\men (\%)} & \makecell[r]{Lower\\bound (\%)} & \makecell[r]{Upper\\bound (\%)} & \makecell[r]{No recorded\\interval (\%)} & \makecell[r]{Weighted\\(\%)} \\
\midrule
Metro, 1 million or more & 205,980 & 42.5 & 40.2 & 45.7 & 5.5 & 42.3 \\
Metro, 250,000 to 1 million & 75,790 & 37.6 & 35.8 & 40.6 & 4.8 & 37.4 \\
Metro, under 250,000 & 27,050 & 34.2 & 33.0 & 36.6 & 3.6 & 34.1 \\
Nonmetro, adjacent to metro & 24,050 & 33.2 & 32.2 & 35.1 & 2.9 & 33.1 \\
Nonmetro, not adjacent to metro & 14,780 & 32.2 & 31.3 & 34.0 & 2.7 & 32.0 \\
Unknown & 170 & 33.3 & 31.9 & 36.1 & 4.2 & 33.5 \\
\bottomrule
\end{tabular}
\par\smallskip
\begin{minipage}{\linewidth}\scriptsize Percentage waiting more than 90 days among treated men meeting cohort steps 1 to 6 (intervals of 0 days excluded). Lower bound: every man without a recorded interval waited 90 days or less; upper bound: every such man waited longer. Weighted: men with a recorded interval, weighted by the inverse of their predicted probability of having one. Numbers of men rounded to the nearest 10. Neither approach addresses selection into recorded treatment.\end{minipage}
\end{table}
